\documentclass[11pt]{article}

% ---------------------------------------------------------------------------
%  A research note on the TTM algorithm, its Tucker-tensor-train (T3)
%  generalization (T3M), the hierarchical-Tucker (HT) view, and the idea of
%  converting between formats to reduce "bubbling".
%
%  Compile with pdflatex (needs: amsmath, amssymb, amsthm, algorithm,
%  algpseudocode, booktabs, hyperref, geometry -- all standard).
%
%  This is an *idea capture* note, not a paper: it records connections we
%  worked out and a hypothesis we are deliberately NOT implementing yet.
% ---------------------------------------------------------------------------

\usepackage[margin=1in]{geometry}
\usepackage{amsmath,amssymb,amsthm}
\usepackage{algorithm}
\usepackage{algpseudocode}
\usepackage{booktabs}
\usepackage[colorlinks=true,linkcolor=blue,citecolor=blue,urlcolor=blue]{hyperref}

\newcommand{\T}{\mathcal{T}}
\newcommand{\A}{\mathcal{A}}
\newcommand{\B}{\mathcal{B}}
\newcommand{\C}{\mathcal{C}}
\newcommand{\R}{\mathbb{R}}
\newcommand{\krrow}{\odot_{\mathrm{r}}}     % row-wise Khatri--Rao
\newcommand{\kron}{\otimes}
\newcommand{\had}{\odot}                     % Hadamard / elementwise product
\newcommand{\rtil}{\tilde{r}}
\newcommand{\ntil}{\tilde{n}}
\DeclareMathOperator{\rank}{rank}
\DeclareMathOperator{\TT}{TT}

\title{\bfseries Elementwise multiplication of tensor trains and Tucker tensor
trains:\\ the TTM algorithm, its T3 generalization, and the hierarchical-Tucker view}
\author{T3Toolbox notes --- N.\ Alger, B.\ Christierson \&\ collaborators}
\date{\today}

\begin{document}
\maketitle

\begin{abstract}
This note records the relationship between three algorithms for the elementwise
(Hadamard) product of tensors stored in tensor-network form: (i) the
\emph{tensor-train multiplication} (TTM) swap algorithm of Michailidis,
Fenton \&\ Kiffner (\href{https://arxiv.org/abs/2410.19747}{arXiv:2410.19747});
(ii) our generalization to Tucker tensor trains (T3), which we call \emph{T3M};
and (iii) the hypothetical generalization to arbitrary hierarchical-Tucker (HT)
trees. We give the product structure, gauge-aware algorithm boxes, and the
complexity. The central conceptual point is that TTM's per-swap truncation is
optimal \emph{only because pure tensor trains have trivial leaf frames}; the
moment leaf frames become nontrivial (the ``Tucker'' in T3, or any HT format),
the leaf compression couples the two operands and forces an
\emph{oversample-then-round} structure. We close with a thought experiment:
converting a TT to a balanced HT to reduce ``bubbling,'' and an analysis of why
it does not pay off for TT/T3 inputs even though the meta-principle is sound.
\end{abstract}

\section{Setup and notation}

A dense order-$d$ tensor $\T\in\R^{N_0\times\cdots\times N_{d-1}}$ has entries
$\T[x_0,\dots,x_{d-1}]$. We write the \emph{Hadamard} (elementwise) product of two
same-shape tensors as $\C=\A\had\B$, i.e.\ $\C[x]=\A[x]\,\B[x]$.

\paragraph{Tensor train (TT).}
$\displaystyle \T[x]=\sum G_0[\,\cdot,x_0,\cdot\,]\cdots G_{d-1}[\,\cdot,x_{d-1},\cdot\,]$,
with cores $G_i\in\R^{r_i\times N_i\times r_{i+1}}$ and boundary bonds
$r_0=r_d=1$. The $r_i$ are the \emph{TT bonds}.

\paragraph{Tucker tensor train (T3).}
A Tucker decomposition whose core $\C$ is itself stored as a TT:
\begin{equation}
\T[x]=\sum_{n_0,\dots,n_{d-1}} \Big(\prod_{i=0}^{d-1} U_i[n_i,x_i]\Big)\,\C[n_0,\dots,n_{d-1}],
\qquad \C=\TT(G),
\label{eq:t3}
\end{equation}
with \emph{Tucker factors} (leaf frames) $U_i\in\R^{n_i\times N_i}$ and central TT
cores $G_i\in\R^{r_i\times n_i\times r_{i+1}}$. The $n_i$ are the \emph{Tucker ranks}.
Plain TT is the special case $n_i=N_i$, $U_i=I$.

\paragraph{Two combination rules.}
For a Hadamard product, the row-wise Khatri--Rao and the Kronecker product appear:
\begin{align}
(U^A\krrow U^B)[(a,b),\,x] &= U^A[a,x]\,U^B[b,x] && \text{(row-wise Khatri--Rao, ``COPY'')},\\
(G^A\kron G^B)[(\alpha,\alpha'),(a,b),(\beta,\beta')] &= G^A[\alpha,a,\beta]\,G^B[\alpha',a',\beta'] && \text{(Kronecker)} .
\end{align}

\paragraph{The product of two T3s.}
Substituting \eqref{eq:t3} for $\A$ and $\B$ into $\C[x]=\A[x]\B[x]$ and grouping
factors mode by mode gives the exact product T3:
\begin{equation}
W_i \;=\; U^A_i\krrow U^B_i \in\R^{(n^A_i n^B_i)\times N_i},
\qquad
G^C_i \;=\; G^A_i\kron G^B_i \in\R^{(r^A_i r^B_i)\times(n^A_i n^B_i)\times(r^A_{i+1}r^B_{i+1})}.
\label{eq:prod}
\end{equation}
\emph{Both} rank families blow up multiplicatively: Tucker rank
$n^A_i n^B_i$ and TT bond $r^A_i r^B_i$. The point of every algorithm below is to
produce \eqref{eq:prod} \emph{already truncated} back to target ranks
$\ntil_i,\rtil_i$, ideally without ever forming the blown-up factors.

\section{The TTM algorithm (pure tensor trains)}
\label{sec:ttm}

For pure TT the leaf frames are trivial, so only the central object matters and
the Hadamard acts \emph{diagonally} on each physical leg: when two same-mode cores
meet, a copy tensor $\delta[x,x',x'']$ (the order-3 diagonal) ties their physical
legs to a single output leg of the \emph{same} size $N_i$. TTM never forms the
$r^2$ bonds. Instead it views $\A\kron\B$ (the order-$2d$ outer product) as one
length-$2d$ TT and merges matching modes by local, gauge-centered, truncating
\emph{swaps}.

\begin{algorithm}[t]
\caption{\textsc{TTM} --- Hadamard product of two tensor trains (swap form)}
\label{alg:ttm}
\begin{algorithmic}[1]
\Require TT cores $\{G^A_i\}$, $\{G^B_i\}$; per-step tolerance / target bond $\rtil$
\State $H \gets [\,G^A_0,\dots,G^A_{d-1},\ \textsc{reverse}(G^B)\,]$ \Comment{length-$2d$ chain $=\A\kron\B$; reversal makes mode $d{-}1$ adjacent}
\State left-orthogonalize the $A$-half, right-orthogonalize the reversed-$B$ half \Comment{mixed-canonical center at the seam}
\For{$i=d-1$ \textbf{downto} $0$}
  \State move the orthogonality center to the right end of the merged block
  \For{each step bringing the mode-$i$ $B$-core toward its $A$-partner}
    \State \textsc{swap} the adjacent pair by a \emph{centered} truncated SVD \Comment{bond stays $O(\rtil)$}
  \EndFor
  \State \textsc{merge} the adjacent matching pair with the copy tensor $\delta$ (diagonal in $x_i$); single bonds in/out
\EndFor
\State \Return the length-$d$ chain $\{G^C_i\}$
\end{algorithmic}
\end{algorithm}

The reversal makes the $\tfrac{d(d-1)}{2}$ swaps minimal, and each merged core is
``bubbled past'' by later $B$-cores. Because every swap and merge is performed in
mixed-canonical gauge (everything surrounding the active core is orthogonal toward
it), each truncated SVD is locally optimal in Frobenius norm. \textbf{For pure TT
this per-step optimality is enough}: the physical leg size $N_i$ is invariant
(the diagonal merge does not compress it), so the only truncation is of the
\emph{bonds}, and centered bond truncation is optimal. Cost $O(d^2\rtil^3)$,
peak memory $O(\rtil^2)$ --- the regime win is $\rtil\gg d$.

\section{Generalization to Tucker tensor trains (T3M)}
\label{sec:t3m}

\subsection{What changes, and what does not}
The skeleton of Algorithm~\ref{alg:ttm} ports verbatim: concatenate, reverse,
the $\tfrac{d(d-1)}{2}$ swap schedule, and the gauge discipline are all
tensor-agnostic. \emph{One} thing changes: the copy-tensor merge becomes the
\textbf{Khatri--Rao merge}. When the mode-$i$ pair meets, the merged physical leg
is the pair $(n^A_i,n^B_i)$ carrying the Tucker factor $W_i=U^A_i\krrow U^B_i$ from
\eqref{eq:prod}, and $W_i$ admits a \emph{new} compression
\begin{equation}
\rank(W_i)\le \min(n^A_i n^B_i,\;N_i),\qquad\text{truncate to } \ntil_i,
\end{equation}
weighted by the (orthogonal) environment exactly as the T3-SVD per-site step does.
This Tucker compression \textbf{has no analogue in pure TT} and is the source of
everything below.

\begin{algorithm}[t]
\caption{\textsc{T3M-swap} --- Hadamard product of two Tucker tensor trains}
\label{alg:t3m}
\begin{algorithmic}[1]
\Require T3s $(U^A,G^A)$, $(U^B,G^B)$; targets $\ntil,\rtil$ (and/or $\mathrm{rtol},\mathrm{atol}$); oversample $k\ge1$
\State down-orthogonalize the Tucker factors of each input \Comment{HT: ``orthogonalize leaves toward the root''}
\State $H \gets [\,G^A_0,\dots,G^A_{d-1},\ \textsc{reverse}(G^B)\,]$ with each $G_i$ carrying $(U_i,\text{mode }i)$
\State left-orthogonalize the $A$-half, right-orthogonalize the reversed-$B$ half \Comment{center at the seam}
\For{$i=d-1$ \textbf{downto} $0$}
  \State move the center to the right end of the merged block
  \For{each step bringing the mode-$i$ $B$-core to its $A$-partner}
    \State centered truncated-SVD \textsc{swap}, intermediate bond capped at $k\,\rtil$ (or tol $/k$) \Comment{memory only}
  \EndFor
  \State \textsc{merge}: $W_i\gets U^A_i\krrow U^B_i$; centered weighted-SVD compress to $\le k\,\ntil$; Tucker factors/modes ride along
\EndFor
\If{$k>1$ \textbf{or} a non-uniform per-position bond cap was requested}
  \State $\{(U^C_i,G^C_i)\}\gets$ \textsc{T3-SVD-round}$(\{(U^C_i,G^C_i)\};\ \ntil,\rtil,\mathrm{rtol},\mathrm{atol})$ \Comment{the decisive truncation}
\EndIf
\State \Return $\{(U^C_i,G^C_i)\}$
\end{algorithmic}
\end{algorithm}

\subsection{The leaf-frame tension}
With the gauge done correctly the no-truncation path of Algorithm~\ref{alg:t3m}
reproduces \eqref{eq:prod} to machine precision. Yet under \emph{aggressive}
truncation it lands at $\approx1.2\times$ the optimal error, whereas
form-then-round is optimal. The cause is structural, not a gauge bug:
\begin{quote}
the Khatri--Rao compression $W_i\!=\!U^A_i\krrow U^B_i$ \textbf{couples} the two
operands' factors, so it can only happen when the two cores \emph{meet} (the merge,
which for the last mode is the very last step). Meanwhile the bonds want to be
truncated \emph{after} all merges: while $a_0$ and $b_0$ are still split across the
chain, the mode-$0$ correlation is carried by the \emph{bonds}; once merged it
becomes internal to the merged core and stops burdening them. Truncating bonds
before the final merge truncates a representation with artificially inflated bonds.
\end{quote}
Pure TT escapes this precisely because its leaf frames are trivial: no leaf
compression, no coupling, no tension. \textbf{The tension is the price of
nontrivial leaf frames.}

\subsection{Oversample, then round}
The principled resolution: use the swaps only to bring cores together with
\emph{loosely} bounded intermediate ranks (a factor $k\ge1$ of oversampling,
purely for memory), and defer the \emph{decisive} truncation to a final
T3-SVD rounding pass --- run after every Tucker leg is already compressed and with
$O(d)$ rather than $O(d^2)$ steps. Empirically $k\!\approx\!3$ recovers the optimal
error; $k\!=\!1$ (no extra rounding) is the maximal-memory-saving, lowest-quality
corner that coincides with the naive ``truncate hard in place'' reading of TTM.
Crucially, in the $\rtil\gg d$ regime where the swap method is used, $\rtil\ll r$, so
$O((k\rtil)^2)$ memory with small $k$ is still far below the $O(r^2)$ full product:
the memory advantage survives. The final round, being a native T3-SVD, also makes
a \emph{per-position} bond-cap sequence honored exactly --- something the swap
process cannot do bond-by-bond, since it does not preserve per-position bond
identity.

\section{The view from hierarchical Tucker (HT)}
\label{sec:ht}

\subsection{TT, Tucker, and T3 as HT instances}
An HT tensor is a \emph{dimension tree} with leaf frames $U_i\in\R^{n_i\times N_i}$
and transfer tensors $B_t\in\R^{n_{t_1}\times n_{t_2}\times n_t}$ at internal nodes
$t$ with children $t_1,t_2$ (root rank $1$). The familiar formats are points in
this space:
\begin{center}
\begin{tabular}{ll}
\toprule
format & HT realization\\
\midrule
Tensor train & linear (caterpillar) tree, \textbf{trivial} leaf frames ($U_i=I$)\\
Tucker & star (one root transfer tensor, $d$ leaf frames)\\
Tucker tensor train (T3) & linear tree, \textbf{nontrivial} leaf frames\\
\bottomrule
\end{tabular}
\end{center}
So T3 is exactly ``TT with the leaf frames switched on'' --- which, by
\S\ref{sec:t3m}, is the whole story behind the tension.

\subsection{Hadamard on trees}
For two HT tensors on the \emph{same} tree, the Hadamard product is, at every node,
\[
\text{leaf frames } \;W_i = U^A_i\krrow U^B_i,
\qquad
\text{transfer tensors } \; B^C_t = B^A_t\kron B^B_t,
\]
i.e.\ the universal ``KR at leaves, Kronecker inside,'' followed by an HT-rounding.
Equivalently (the TTM viewpoint), $\A\kron\B$ is itself a valid HT tensor on the
two trees joined at a new root, and the Hadamard is recovered by \emph{merging each
pair of matching leaves} (the copy tensor / KR) while truncating in mixed-canonical
gauge. TTM/T3M are the \emph{linear-tree} instances of this; a general tree replaces
``adjacent swaps'' with SVD-based subtree/gauge moves.

This lens retro-justifies several T3M choices as canonical HT steps:
down-orthogonalizing the Tucker factors $=$ orthogonalizing leaves toward the root;
the center-at-the-truncated-node gauge $=$ HT mixed-canonical form; the final
T3-SVD cleanup $=$ ``round with the format's native rounder.''

\subsection{Bubbling and tree balance}
In a linear tree the matching leaves of $\A\kron\B$ start $\sim d$ apart, and
rearranging them to be siblings costs $O(d^2)$ adjacent swaps (the ``bubbling'';
it is the inversion count of the permutation
$[0,\dots,d{-}1,\,d{-}1,\dots,0]\mapsto[0,0,1,1,\dots]$). A \emph{balanced} tree has
diameter $O(\log d)$, so the rearrangement count drops to $O(d\log d)$. Thus the
$O(d^2)$ in T3M is a \emph{linear-tree artifact}, intrinsic to T3 (which is linear
by definition), consistent with positioning the swap method for small $d$, large
$\rtil$.

\section{Thought experiment: convert, multiply, convert back}
\label{sec:convert}

If a balanced tree needs fewer bubbles, why not \emph{convert} a TT/T3 into a
balanced HT, take the Hadamard there, and convert back? We record the idea and the
reason it does \emph{not} pay off for TT/T3 inputs --- while noting the
meta-principle is nonetheless sound.

\paragraph{Rank inflation under TT$\to$HT.}
A contiguous \emph{middle} block $[a,b]$ of a TT connects to the rest through
\emph{two} bonds, so its HT separation rank obeys
$\rank_{[a,b]}\le r_a\,r_{b+1}\approx r^2$ (prefix/suffix blocks keep rank
$\approx r$; the middle nodes inflate). Since the separation rank of a Hadamard
product is sub-multiplicative,
\[
\rank_{[a,b]}(\A\had\B)\ \le\ \rank_{[a,b]}(\A)\,\rank_{[a,b]}(\B)\ \approx\ r^2\cdot r^2 = r^4 .
\]
So the product is $\approx r^2$-compact as a \textbf{TT} but $\approx r^4$-compact
as a \textbf{balanced HT}: converting moves you into a representation where the very
object being computed is quadratically larger.

\paragraph{Cost comparison (working bond $\rtil$).}
\begin{center}
\begin{tabular}{lll}
\toprule
 & rearrangement count & per-operation cost \\
\midrule
linear / TT swap & $O(d^2)$ & $O(\rtil^3)$ \\
balanced HT & $O(d\log d)$ & $O((\rtil^2)^3)=O(\rtil^6)$ \\
\bottomrule
\end{tabular}
\end{center}
Hence
\[
\frac{\text{balanced}}{\text{linear}}\ \approx\ \frac{\log d}{d}\,\rtil^3 .
\]
In the $\rtil\gg d$ regime --- exactly where bubbling matters and the swap method is
used --- $\rtil^3\gg d$, so the balanced route is \emph{worse}; in the $d\gg\rtil$
regime the in-place fused method already avoids bubbling, so one would not swap at
all. There is thus no regime where convert-then-convert-back clearly beats the best
existing method for TT/T3 inputs.

\paragraph{What survives.}
The \emph{meta-principle} --- compute in whatever network topology makes the
operation cheapest, then convert back --- is legitimate and used elsewhere in
tensor-network practice. It simply requires the alternative topology to be
genuinely more compact \emph{for this data and this operation}, and balanced HT is
not, for TT-Hadamard. A partial / data-adaptive rebalancing (cheaper than full
balancing, avoiding the worst middle-block inflation) is a genuine open direction,
but is out of scope here.

\section{Summary}

\begin{center}
\begin{tabular}{lllp{4.6cm}}
\toprule
method & compute & peak memory & character \\
\midrule
(a) form-then-round & $O(d\,r^6)$ & $O(d\,r^4 n^2)$ & forms full product; embarrassingly parallel; the optimal-quality oracle \\
(b) in-place fused  & $O(d\,r^4)$ & $O(\rtil\,n^2 r^2)$ & fused L$\to$R sweep; the default for large $d$ \\
(c) swap (T3M)      & $O(d^2 r^3)$ & $O(\rtil^2)$ & this note; for $r\gg d$; needs oversample$+$round \\
\midrule
hypothetical balanced HT & $O(d\log d\,\rtil^6)$ & $O(\rtil^4)$ & fewer bubbles, but rank-inflated; not a win for TT/T3 \\
\bottomrule
\end{tabular}
\end{center}

\paragraph{Take-aways.}
(1) TTM's swap skeleton is the faithful generalization to T3; only the merge changes
(copy tensor $\to$ Khatri--Rao). (2) Per-swap truncation is optimal \emph{only} for
trivial leaf frames; nontrivial leaf frames (T3, HT) introduce a leaf-compression
coupling that forces \emph{oversample-then-round}. (3) T3 is the linear-tree,
leaf-bearing instance of an HT Hadamard algorithm; its $O(d^2)$ bubbling is a
linear-tree artifact. (4) Converting to a balanced HT trades bubble count for rank
inflation and does not help for TT/T3 inputs, though the underlying
topology-choice principle is sound.

\paragraph{References.}
P.\ Michailidis, C.\ Fenton, M.\ Kiffner, \emph{Tensor Train Multiplication},
\href{https://arxiv.org/abs/2410.19747}{arXiv:2410.19747} (TTM).\\
N.\ Alger, B.\ Christierson, P.\ Chen, O.\ Ghattas, \emph{Tucker Tensor Train
Taylor Series}, arXiv:2603.21141 (T4S; T3 format, T3-SVD).

\end{document}
